The ideal generated by @@M0@@ is a @@M1@@-ideal. To see this, first observe that the relations @@M2@@ for @@M3@@ imply that the @@M4@@ for @@M5@@ are multiples of @@M6@@. Working in @@M7@@, the image of @@M8@@ is a polynomial of degree @@M9@@, and a universal argument then shows that its @@M10@@ for @@M11@@ equal @@M12@@ (cf. V 4.8). Since in @@M13@@ we have @@M14@@, we deduce that the @@M15@@ for @@M16@@ are multiples of @@M17@@. By the relations V 2.4.2, this implies that the ideal generated by @@M18@@ is a @@M19@@-ideal.

Put then @@M20@@. We denote the natural homomorphism @@M21@@ by @@M22@@; it is a @@M23@@-homomorphism. Furthermore, @@M24@@ is a free @@M25@@-module with basis the elements @@M26@@. The relation @@M27@@ and the invertibility hypothesis on @@M28@@ imply that @@M29@@ is invertible. Consequently, @@M30@@ has as a basis over @@M31@@ the elements @@M32@@; we then define a map @@M33@@ by assigning to every element @@M34@@ of @@M35@@ the element @@M36@@ of @@M37@@. Clearly the map thus defined is an additive, @@M38@@-linear homomorphism, satisfying @@M39@@.

Next we endow @@M40@@ with a @@M41@@-ring structure, putting:
\sgaStableTarget{sga6.e07.eq.4-7-1}\sgaStableTarget{sga6.e07.eq.4-7-2}@@M42@@
In these formulas, @@M43@@ is regarded as a @@M44@@-algebra through the homomorphism @@M45@@.
Next, for @@M46@@ and @@M47@@, put:
\sgaStableTarget{sga6.e07.eq.4-7-3}\sgaStableTarget{sga6.e07.eq.4-7-4}@@M48@@
